% ============================================================
% PT Article M2 — version française
% Miroir FR de m2_geometry.tex (canonique EN)
% Préambule : charge pt-common-fr.sty (shared/) qui contient macros,
% boîtes, environnements de théorème FR, patch babel/cleveref, etc.
% ============================================================

\documentclass[12pt,a4paper]{article}
\IfFileExists{pt-common-fr.sty}{\usepackage{pt-common-fr}}{\usepackage{../shared/pt-common-fr}}
\usepackage[margin=2.5cm]{geometry}
\usepackage{setspace}
\usepackage{tikz-3dplot}
\onehalfspacing

\title{\textbf{La théorie de la persistance~II~:\\
géométrie de l'information et holonomie}}
\author{Yan Senez\\\texttt{yan.senez@protonmail.com}}
\date{Avril 2026}

\begin{document}
\maketitle

\begin{abstract}
Ceci est le second de cinq articles présentant les fondations
mathématiques de la théorie de la persistance (PT). En s'appuyant sur la
dynamique du crible, l'unicité, les lois de conservation et le théorème
fondamental des sauts établis dans l'article
compagnon~\cite{companion_M1}, nous développons deux piliers de la
théorie~: la géométrie canonique de l'information de l'espace d'états du
crible, et la structure d'holonomie qui émerge des matrices de transfert
du crible.

\textbf{Géométrie de l'information.}
La divergence de KL $\DKL$ est l'unique $f$-divergence compatible avec
la structure produit~TCR de l'espace d'états du crible
(théorème~G1, par réduction à Shore--Johnson~\cite{ShoreJohnson1980},
avec 7/7 tests d'élimination~; une preuve autonome via TCR utilisant la
caractérisation de Csisz\'ar~\cite{Csiszar1967} est également donnée).
La métrique de Fisher $\Fisher$ émerge comme seconde variation de
$\DKL$ (théorème~G2, classique) et est l'unique métrique riemannienne
monotone sur $\Delta_2$ (réduction à \v{C}encov~\cite{Cencov1982},
théorème~G3, 7/7 tests d'élimination). Nous exhibons la métrique de
Fisher sur les simplexes binaire et ternaire, montrons qu'elle est
imposée par la structure d'expanseur optimal du graphe de transition
mod-3, et la relions à la dualité addition/multiplication des entiers.

\textbf{Holonomie.}
L'angle d'holonomie $\theta_p$ satisfait l'identité algébrique exacte
$\sinsp{p} = \delta_p(2 - \delta_p)$, dérivée par trois voies
indépendantes (géométrique, spectrale, information-géométrique). Les
dimensions anomales $\gamma_p = -d\ln\sinsp{p}/d\ln\mu$ sont strictement
décroissantes en $p$~; le critère de premier actif $\gamma_p > s = 1/2$
sélectionne exactement $\{3, 5, 7\}$ au point fixe $\mustar = 15$, ce
qui donne $N_c = 3$. Le coefficient de Casimir
$C_F = (N_c^2 - 1)/(2N_c) = 4/3$ en est une conséquence, et non une
entrée. Un contrôle algébrique indépendant donne
$(N_c - 3)(N_c + 1) = 0$. Les premiers d'écho ($p \geq 11$) contribuent
perturbativement comme corrections de polarisation du vide.

Tous les résultats découlent de $s = 1/2$ et de la structure du crible
établie dans~\cite{companion_M1}. Aucune physique n'est invoquée~;
aucun paramètre n'est ajusté.
\end{abstract}

\tableofcontents
\newpage

\section{Introduction}
\label{sec:intro}

Cet article est le second volet (M2) d'une série de cinq articles
développant les fondations mathématiques de la théorie de la
persistance (PT). La série est organisée comme suit~:
\begin{description}[nosep]
\item[M1~\cite{companion_M1}~:] Le crible comme champ dynamique,
  unicité d'Ératosthène, lois de conservation et théorème fondamental
  des sauts.
\item[M2 (cet article)~:] Géométrie de l'information, holonomie et
  échelle d'angle interne.
\item[M3~\cite{companion_M3}~:] Convergence, récurrences exactes et
  point fixe $\mustar = 15$.
\item[M4~\cite{companion_M4}~:] Structures mathématiques étendues
  (mécanique complexe, graphes spectraux, symétrie).
\item[M5~\cite{companion_M5}~:] Le pont de l'arithmétique vers la
  physique.
\end{description}
\noindent
Un compagnon séparé~\cite{companion_physics} développe les prédictions
physiques détaillées (programme de Feynman, 43~observables)~; M5
fournit le pont structurel, tandis que~\cite{companion_physics} fournit
la validation quantitative.

L'article compagnon~\cite{companion_M1} a établi la couche fondatrice.
Nous rappelons les résultats essentiels~:
\begin{itemize}[nosep]
\item \textbf{T1 (transitions interdites)~:} la matrice de transfert
  $T_3 = \mathrm{antidiag}(1,1)$ admet une unique distribution
  stationnaire $(1/2, 1/2)$, ce qui donne $s = 1/2$ sans paramètre
  libre.
\item \textbf{T6 (unicité)~:} trois preuves indépendantes qu'Ératosthène
  est l'unique crible auto-cohérent.
\item \textbf{L0 (conservation)~:} la distribution des sauts à entropie
  maximale est géométrique~; l'exposant de conservation est
  $\alpha_{\rm cons} = s^2 = 1/4$.
\item \textbf{TFS~:} l'identité exacte $\log_2 m = \DKL + H$ décompose
  la capacité totale en structure et bruit.
\item \textbf{Connexion de jauge TCR~:} le théorème chinois des restes
  dote l'espace d'états du crible $\Delta_m$ d'une structure produit
  $\Delta_m \cong \prod_{p|m} \Delta_p$.
\end{itemize}

Le présent article développe deux thèmes.

\textbf{Géométrie de l'information (section~\ref{sec:geometry}).}
Nous prouvons que la divergence de KL $\DKL$ et la métrique de Fisher
$\Fisher$ sont l'\emph{unique} divergence et l'unique métrique
compatibles avec les symétries du crible. Les preuves sont des
réductions à des théorèmes
classiques~--- Shore--Johnson~\cite{ShoreJohnson1980} pour $\DKL$ et
\v{C}encov~\cite{Cencov1982} pour $\Fisher$~--- après vérification que
l'espace d'états du crible satisfait leurs prémisses. Une preuve
autonome de G1 via la caractérisation de Csisz\'ar des $f$-divergences
additives est également donnée. Cette section est entièrement
mathématique et n'invoque aucune physique.

\textbf{Holonomie (section~\ref{sec:holonomy}).}
La matrice de transfert du crible à chaque premier $p$ engendre un
angle de rotation $\theta_p$ sur l'espace des résidus
$\mathbb{Z}/p\mathbb{Z}$. L'identité d'holonomie
$\sinsp{p} = \delta_p(2 - \delta_p)$ est une tautologie algébrique,
vérifiée par trois voies indépendantes. Les dimensions anomales
$\gamma_p$ sont strictement décroissantes, et le critère
$\gamma_p > s = 1/2$ sélectionne exactement trois premiers actifs
$\{3, 5, 7\}$ à $\mustar = 15$, ce qui donne $N_c = 3$ et
$C_F = 4/3$. L'identification physique de ces quantités arithmétiques
avec les couplages de jauge et les couleurs QCD est un énoncé de pont
reporté à~\cite{companion_M5}.

\medskip\noindent\textbf{Aiguillages amont.}
L'article compagnon~\cite{companion_M3} prouve la convergence des
récurrences du crible et établit $\mustar = 15$ comme unique point
fixe non trivial. Les structures mathématiques étendues (mécanique
complexe, graphes spectraux et extensions de symétrie) sont développées
dans~\cite{companion_M4}. Le pont de l'arithmétique vers la
physique~--- où $\sinsp{p}$ devient une constante de couplage et
$N_c = 3$ devient le nombre de couleurs QCD~--- est l'objet
de~\cite{companion_M5}.

\section{Géométrie canonique de l'information}
\label{sec:geometry}

\epigraph{La géométrie est la branche des mathématiques qui étudie les propriétés
des figures. La géométrie de l'information est la branche qui étudie
les propriétés des variétés statistiques.}{Shun-ichi Amari}

% ---- Status Box (R1) ----------------------------------------
\begin{statusbox}
\textbf{OBJECTIF~:}
  Prouver que les deux divergences canoniques sur l'espace d'états du
  crible~--- la divergence de KL $\DKL$ et la métrique de Fisher
  $\Fisher$~--- sont les \emph{uniques} choix compatibles avec les
  symétries du crible.

\textbf{ENTRÉES~:}
  TCR, connexion de jauge, espace d'états $\Delta_m$
  (\cite{companion_M1}).
  TFS, $\DKL$ déjà utilisée (\cite{companion_M1}).

\textbf{RÉSULTATS PRINCIPAUX~:}
  G1~: $\DKL$ est l'unique $f$-divergence invariante sous TCR
  (réduction à Shore--Johnson, 7/7).
  G2~: Seconde variation de $\DKL$ = métrique de Fisher (classique,
  Amari 1985).
  G3~: La métrique de Fisher est l'unique métrique riemannienne
  monotone sur $\Delta_2$ (réduction à \v Cencov, 7/7).

\textbf{STATUT~:}
  \THMTAG~(réduction~: G1 via Shore--Johnson, G3 via \v Cencov).
  Ce chapitre ne prouve PAS Shore--Johnson ni \v Cencov de façon
  autonome. Il vérifie que l'espace d'états du crible satisfait
  leurs prémisses.

\textbf{NON PROUVÉ~:}
  Shore--Johnson (1980) et \v Cencov (1982)~: importés comme théorèmes
  classiques nommés. Utilisation physique de la métrique de Fisher
  (article compagnon~\cite{companion_physics}).
\end{statusbox}

% ============================================================
\subsection{Intuition~: pourquoi ces deux objets~?}
\label{sec:ch5_intuition}

Le théorème fondamental des sauts (voir~\cite{companion_M1}) utilise
$\DKL$ comme mesure de la structure du crible. L'article
compagnon~\cite{companion_physics} utilisera la métrique de Fisher
pour dériver la relativité générale. Mais pourquoi $\DKL$ et Fisher~?
Parmi l'infinité de divergences et de métriques possibles sur les
distributions de probabilité, qu'est-ce qui distingue ces deux-là~?

La réponse pour $\DKL$~: c'est l'unique $f$-divergence additive sous la
factorisation TCR dont jouit le crible. Toute autre $f$-divergence
(Hellinger, $\chi^2$, variation totale) viole cette invariance
(vérifié~: 7/7 tests).

La réponse pour Fisher~: c'est l'unique métrique riemannienne sur les
simplexes de probabilité qui se contracte sous les applications de
Markov. Puisque la matrice de transfert du crible $T_m$ est une
application de Markov, toute distance information-géométrique doit se
contracter sous celle-ci. Fisher est l'unique métrique ayant cette
propriété.

Les deux résultats sont des réductions à des théorèmes classiques
(Shore--Johnson~\cite{ShoreJohnson1980},
\v{C}encov~\cite{Cencov1982}). La contribution de PT est de vérifier
que l'espace d'états du crible satisfait les prémisses.

% ============================================================
\subsection{L'espace d'états du crible}
\label{sec:ch5_state_space}

\begin{definition}[Simplexe de probabilité $\Delta_m$]
\label{def:simplex}
L'\emph{espace d'états du crible} au module $m$ est le simplexe de
probabilité~:
\[
  \Delta_m = \Bigl\{ P = (p_0, \ldots, p_{m-1}) \in \mathbb{R}^m
             : p_r \geq 0,\; \sum_r p_r = 1 \Bigr\}.
\]
La distribution du crible $P_m \in \Delta_m$ est le vecteur des
fréquences des résidus de premiers modulo $m$.
\end{definition}

Le crible agit sur $\Delta_m$ via la factorisation TCR~:
pour $m = q_1 \cdots q_k$ (sans facteur carré), l'isomorphisme TCR donne
\[
  \Delta_m \cong \Delta_{q_1} \times \cdots \times \Delta_{q_k},
\]
de sorte que l'espace d'états est un produit cartésien de simplexes
plus petits. Tout objet géométrique (divergence, métrique) sur
$\Delta_m$ doit respecter cette structure produit.
La figure~\ref{fig:fisher_simplex} illustre la métrique de Fisher sur
le simplexe binaire $\Delta_2$, et la
figure~\ref{fig:fisher_3d} montre l'extension au simplexe ternaire
$\Delta_3$.

\begin{figure}[htbp]
\centering
\vspace{0.5cm}
\begin{tikzpicture}[scale=5.0, >=stealth]
  % Euclidean baseline
  \draw[gray!30, thick] (0,0) -- (1,0);
  \node[below, font=\small, gray] at (0,-0.02) {$0$};
  \node[below, font=\small, gray] at (1,-0.02) {$1$};
  \node[below, font=\small, gray] at (0.5,-0.02) {$\tfrac{1}{2}$};
  \node[font=\scriptsize, gray] at (0.5,-0.09) {(euclidienne)};
  % Fisher arc
  \draw[thmblue, very thick, domain=0:180, samples=80]
    plot ({0.5 + 0.45*cos(\x)}, {0.45*sin(\x)});
  \node[above, font=\small, thmblue] at (0.5, 0.53) {arc de Fisher};
  % Grid lines showing metric stretching near boundaries
  \foreach \p in {0.05, 0.1, 0.2, 0.3, 0.4, 0.5, 0.6, 0.7, 0.8, 0.9, 0.95} {
    \pgfmathsetmacro{\ang}{asin(sqrt(\p))*2}
    \pgfmathsetmacro{\fx}{0.5 + 0.45*cos(\ang)}
    \pgfmathsetmacro{\fy}{0.45*sin(\ang)}
    \draw[thmblue!40, thin] (\p, 0) -- (\fx, \fy);
  }
  % Stationary point p = 1/2 — leader line below Fisher arc label
  \fill[condorange] (0.5, 0.45) circle (0.7pt);
  \draw[condorange, thin] (0.5, 0.45) -- (0.68, 0.48);
  \node[right, font=\small, condorange] at (0.68, 0.48) {$s = \tfrac{1}{2}$};
  % Stretching near boundaries
  \draw[<->, thin, predred] (0.02, -0.12) -- (0.12, -0.12);
  \node[below, font=\scriptsize, predred] at (0.07, -0.13) {$ds \to \infty$};
  \draw[<->, thin, predred] (0.88, -0.12) -- (0.98, -0.12);
  \node[below, font=\scriptsize, predred] at (0.93, -0.13) {$ds \to \infty$};
  % Metric formula
  \node[font=\normalsize] at (0.5, -0.17)
    {$ds^2 = \dfrac{dp^2}{p\,(1-p)}$};
  % axis label
  \node[below, font=\small] at (0.5, -0.28) {probabilité $p = \pi_1$};
\end{tikzpicture}
\caption{Métrique de Fisher sur le simplexe binaire $\Delta_2$.
  L'intervalle euclidien $[0,1]$ (gris) est «~courbé~» par la métrique
  de Fisher $ds^2 = dp^2/[p(1-p)]$ en un demi-cercle (arc bleu).
  Les lignes auxiliaires relient des valeurs de probabilité égales~:
  elles se resserrent près des bords ($p \to 0$ ou $p \to 1$),
  signalant la divergence de la métrique.
  Le point stationnaire $s = 1/2$ (orange) se trouve au milieu
  géodésique. C'est l'\emph{unique} métrique monotone sur $\Delta_2$
  (\v Cencov).}
\label{fig:fisher_simplex}
\end{figure}

\begin{figure}[htbp]
\centering
\tdplotsetmaincoords{68}{125}
\begin{tikzpicture}[tdplot_main_coords, scale=4.0, >=stealth]
  % Ternary simplex Delta_3: vertices
  \coordinate (P1) at (1,0,0);
  \coordinate (P2) at (0,1,0);
  \coordinate (P3) at (0,0,1);
  % Simplex edges
  \draw[gray!50, thick] (P1) -- (P2) -- (P3) -- cycle;
  % Fill simplex face
  \fill[thmblue!8, opacity=0.4] (P1) -- (P2) -- (P3) -- cycle;
  % Grid lines on the simplex (explicit coordinates, no calc)
  \foreach \i in {1,...,8} {
    \pgfmathsetmacro{\t}{\i/9}
    \pgfmathsetmacro{\omt}{1-\t}
    % Lines parallel to P1-P2 edge
    \draw[thmblue!30, very thin]
      ({\omt}, 0, {\t}) -- (0, {\omt}, {\t});
    % Lines parallel to P2-P3 edge
    \draw[thmblue!30, very thin]
      ({\t}, {\omt}, 0) -- ({\t}, 0, {\omt});
    % Lines parallel to P1-P3 edge
    \draw[thmblue!30, very thin]
      (0, {\t}, {\omt}) -- ({\omt}, {\t}, 0);
  }
  % Centroid = (1/3, 1/3, 1/3)
  \fill[condorange] (0.333, 0.333, 0.333) circle (0.8pt);
  \draw[condorange, thin] (0.333, 0.333, 0.333) -- (0.6, 0.6, 0.5);
  \node[right, font=\small, condorange] at (0.6, 0.6, 0.5)
    {$p_i = \tfrac{1}{3}$ ($s = \tfrac{1}{2}$)};
  % Vertex labels
  \node[below right, font=\small] at (P1) {$p_1 = 1$};
  \node[below left, font=\small] at (P2) {$p_2 = 1$};
  \node[above, font=\small] at (P3) {$p_3 = 1$};
  % Edge labels (divergence)
  \node[font=\scriptsize, predred] at (0.5, 0.5, 0) {$ds \to \infty$};
  \node[font=\scriptsize, predred] at (0, 0.5, 0.5) {$ds \to \infty$};
  \node[font=\scriptsize, predred] at (0.5, 0, 0.5) {$ds \to \infty$};
  % Metric formula — offset left to avoid overlapping triangle
  \node[font=\small, align=center, gray] at (-0.6, -0.2, 0.5)
    {$ds^2 = \sum_i \dfrac{dp_i^2}{p_i}$};
\end{tikzpicture}
\caption{Métrique de Fisher sur le simplexe ternaire $\Delta_3$.
  Les trois sommets (états purs, $p_i = 1$) se trouvent à distance de
  Fisher infinie de tout point intérieur~--- la métrique diverge près
  des arêtes. Les lignes auxiliaires montrent comment des pas
  euclidiens égaux correspondent à des longueurs de Fisher inégales
  (resserrement près des bords). Le centroïde $p_i = 1/3$ (orange,
  correspondant à $s = 1/2$ dans la projection binaire) est le milieu
  géodésique.}
\label{fig:fisher_3d}
\end{figure}

% ============================================================
\subsection{G1~: unicité de $\DKL$ via Shore--Johnson}
\label{sec:G1}

\paragraph{Le théorème de Shore--Johnson (1980, import externe).}
Sous les axiomes suivants sur une divergence $D(\cdot\|\cdot)$~:
\begin{enumerate}[label=(SJ\arabic*)]
\item \textbf{Cohérence}~: $\arg\min_Q D(Q\|P)$ sous contraintes est
  indépendant de l'ordre dans lequel les contraintes sont imposées.
\item \textbf{Invariance}~: $D$ est invariant sous permutation des
  étiquettes.
\item \textbf{Indépendance des systèmes}~: pour des systèmes
  indépendants, $D$ est additive.
\item \textbf{Indépendance des sous-ensembles}~: une mise à jour sur
  une contrainte qui ne concerne qu'un sous-ensemble ne modifie que la
  distribution de ce sous-ensemble.
\item \textbf{Échelle}~: $D$ se transforme de manière cohérente sous
  réétiquetage.
\end{enumerate}
Shore et Johnson~\cite{ShoreJohnson1980} ont prouvé que l'unique
divergence satisfaisant SJ1--SJ5 et produisant le principe d'entropie
maximale est $\DKL$.

\mTHM
\begin{theoreme}[Unicité de la divergence de KL (G1)]
\label{thm:G1}\leanPartial{PT.Information.T6cG1Autonomous.klDivergence\_additive}
Sur l'espace d'états du crible $\Delta_m$ muni de la structure produit
TCR~: $\DKL$ est l'unique $f$-divergence satisfaisant les axiomes
SJ1--SJ5 de Shore--Johnson et invariante sous la factorisation TCR.

\begin{proof}
\textbf{Vérification des prémisses.} L'espace d'états du crible
$\Delta_m$ muni de la structure produit
$\Delta_m \cong \prod_{p|m} \Delta_p$ satisfait~:
\begin{itemize}
\item SJ3 (indépendance des systèmes)~: la factorisation TCR donne
  $\DKL(\prod P_p \| \prod Q_p) = \sum_p \DKL(P_p \| Q_p)$.
  L'indépendance des composantes TCR implique l'additivité.
  (Vérifié~: 7/7 tests.)
\item SJ4 (indépendance des sous-ensembles)~: une mise à jour des
  résidus modulo $p$ ne modifie que la composante~$p$ dans la
  décomposition TCR. (Vérifié~: 7/7 tests.)
\item SJ1, SJ2, SJ5~: propriétés standard de cohérence et d'invariance
  satisfaites par la référence uniforme $U_m$ et la distribution
  empirique $P_m$.
\end{itemize}

\textbf{Application de Shore--Johnson.} Puisque les prémisses
SJ1--SJ5 sont vérifiées sur $\Delta_m$, le théorème de Shore--Johnson
donne~: $\DKL$ est l'unique divergence satisfaisant ces axiomes.
(Import externe~: Shore--Johnson 1980, publié dans IEEE
Trans.\ Inf.\ Theory.)

\textbf{Élimination des alternatives.}
\begin{itemize}
\item Divergence $\chi^2$~: viole SJ3 (non additive sous TCR, rapport
  de différences carrées non additif).
\item Distance de Hellinger~: viole SJ3 (la racine carrée brise
  l'additivité).
\item Variation totale~: viole SJ4 (mettre à jour une coordonnée
  affecte globalement la norme~$\ell_1$).
\item Divergences de R\'enyi $D_\alpha$ ($\alpha \neq 1$)~: violent SJ1
  (incohérence des mises à jour séquentielles pour $\alpha \neq 1$).
\end{itemize}
\end{proof}
\end{theoreme}

\begin{remarque}[Portée de G1]
G1 est une \emph{réduction} à Shore--Johnson, et non une preuve
autonome. La contribution PT~: vérifier que $\Delta_m$ muni de la
structure TCR satisfait SJ1--SJ5. La conclusion d'unicité est importée.
\end{remarque}

\begin{theoreme}[Unicité de KL autonome via TCR (G1$'$)]
\label{thm:G1_autonomous}\leanPartial{PT.Information.T6cG1Autonomous.fDivergence\_linear\_term\_vanishes}
Sur l'espace d'états du crible $\Delta_m$ muni de la structure produit
TCR, $\DKL$ est l'unique $f$-divergence satisfaisant l'additivité sous
sous-systèmes indépendants.

\begin{proof}
\textbf{Étape 1 (TCR $\Rightarrow$ structure produit).}
Par le théorème chinois des restes (voir théorème~3.10
de~\cite{companion_M1}), pour $m = q_1 \cdots q_k$ (sans facteur
carré)~:
$\Delta_m \cong \Delta_{q_1} \times \cdots \times \Delta_{q_k}$.
Toute divergence $D$ sur $\Delta_m$ doit respecter cette factorisation.

\textbf{Étape 2 (caractérisation de l'additivité).}
Une $f$-divergence $D_f(P\|Q) = \sum_i q_i\,f(p_i/q_i)$ est additive
sur les distributions produits ($P = P_1 \otimes P_2$,
$Q = Q_1 \otimes Q_2$
$\Rightarrow D_f(P\|Q) = D_f(P_1\|Q_1) + D_f(P_2\|Q_2)$) si et
seulement si $f$ est de la forme
$f(t) = c_1\,t\ln t + c_2(t-1)$ avec des constantes $c_1 > 0$,
$c_2 \in \mathbb{R}$ (Csisz\'ar~\cite{Csiszar1967}).
Puisque $\sum_i q_i(p_i/q_i - 1) = 0$ sur les distributions
normalisées, le terme en $c_2$ s'annule identiquement. L'unique
$f$-divergence additive (à un réétalonnage positif près) est donc
$D_f(P\|Q) = \sum_i p_i \ln(p_i/q_i) = \DKL(P\|Q)$.

\textbf{Étape 3 (unicité).}
Puisque TCR impose la structure produit à chaque niveau du crible, et
puisque $f(t) = t\ln t$ est l'unique générateur produisant une
divergence additive, $\DKL$ est l'unique $f$-divergence additive sur
$\Delta_m$. Aucun import de Shore--Johnson n'est nécessaire~; la
preuve n'utilise que TCR, la définition des $f$-divergences et la
caractérisation de Csisz\'ar des générateurs additifs (analyse convexe
élémentaire).

\textbf{Élimination.}
\begin{itemize}[nosep]
\item $\chi^2$~: $f(t) = (t-1)^2$. Pas de la forme
  $c_1 t\ln t + c_2(t-1)$. Numérique~:
  $D_{\chi^2}(P\otimes P' \| Q\otimes Q') \neq D_{\chi^2}(P\|Q) + D_{\chi^2}(P'\|Q')$ (vérifié).
\item Hellinger~: $f(t) = (\sqrt{t}-1)^2$. Pas de la forme requise.
  Échoue à l'additivité (vérifié).
\item Variation totale~: $f(t) = |t-1|$. Non différentiable en $t=1$.
  Échoue à l'additivité (vérifié).
\end{itemize}

\noindent
Vérifié~: \texttt{scripts/ch05\_geometry/test\_G1\_autonomous.py}
(12/12~PASS).
\end{proof}
\end{theoreme}

\begin{remarque}[Deux preuves indépendantes de G1]
Le théorème~\ref{thm:G1} (réduction à Shore--Johnson) et le
théorème~\ref{thm:G1_autonomous} (TCR autonome) constituent deux
preuves indépendantes de l'unicité de $\DKL$. La première importe
Shore--Johnson~; la seconde n'utilise que la caractérisation de
Csisz\'ar.
\end{remarque}

% ============================================================
\subsection{G2~: la métrique de Fisher comme seconde variation de $\DKL$}
\label{sec:G2}

\mTHM
\begin{theoreme}[Émergence de la métrique de Fisher (G2)]
\label{thm:G2}\leanProved{PT.Information.G2FisherEmergence.hasDerivAt\_klTerm}
Le développement de Taylor au second ordre de $\DKL$ autour de la
distribution uniforme $U_m$ donne la métrique d'information de Fisher~:
\[
  \DKL(P + \epsilon\, v \| P) = \frac{\epsilon^2}{2} \sum_{r,s}
  g_{rs}^F(P)\, v_r\, v_s + O(\epsilon^3),
\]
où $g_{rs}^F(P) = \delta_{rs}/p_r$ est la métrique de Fisher sur
$\Delta_m$.

\begin{proof}
Développement de Taylor de $\DKL(Q\|P) = \sum_r q_r \log(q_r/p_r)$
autour de $Q = P + \epsilon v$~:
\[
  \DKL(P+\epsilon v \| P) = \frac{\epsilon^2}{2} \sum_r \frac{v_r^2}{p_r}
  + O(\epsilon^3).
\]
La hessienne
$\partial^2 \DKL / \partial q_r \partial q_s |_{Q=P} = \delta_{rs}/p_r$
est exactement la métrique de Fisher (résultat classique~:
Amari~\cite{Amari1985}, théorème~2.1).
\end{proof}
\end{theoreme}

% ============================================================
\subsection{G3~: unicité de la métrique de Fisher via \v Cencov}
\label{sec:G3}

\paragraph{Le théorème de \v Cencov (1982, import externe).}
Soit $\mathcal{M}_n$ la variété de toutes les distributions de
probabilité sur $n$ issues. Une métrique riemannienne $g$ sur
$\mathcal{M}_n$ est \emph{monotone} si, pour toute application de
Markov $T$ (matrice stochastique)~:
\[
  g_{T(P)}(T_* v, T_* v) \leq g_P(v, v).
\]
\v{C}encov~\cite{Cencov1982} a prouvé~: la métrique de Fisher est
l'unique métrique riemannienne monotone sur $\mathcal{M}_n$ (à une
constante positive près).

\mTHM
\begin{theoreme}[Unicité de la métrique de Fisher (G3)]
\label{thm:G3}\leanExternal{PT.Information.G3FisherUniqueness.G3\_fisher\_unique}
Sur l'espace d'états du crible $\Delta_2 = \{(p, 1-p) : p \in (0,1)\}$,
la métrique de Fisher $g^F$ est l'unique métrique riemannienne
monotone sous toutes les applications de Markov du crible (y compris
$T_m$ et ses marginales).

\begin{proof}
\textbf{Vérification que $T_m$ est une application de Markov.}
La matrice de transfert $T_m$ est stochastique en lignes (les lignes
somment à~1, entrées $\geq 0$ par définition). Donc $T_m$ est une
application de Markov.

\textbf{Vérification de la prémisse $\Delta_2$.} L'espace d'états du
crible à $m = 3$ est
$\Delta_2 = \{(\pi_1, \pi_2) : \pi_1 + \pi_2 = 1\}$
(deux classes de résidus non nulles, depuis T1~\cite{companion_M1}).
La métrique de Fisher sur $\Delta_2$ est $ds^2 = dp^2/(p(1-p))$, la
métrique de longueur d'arc sur $[0,1]$.

\textbf{Monotonicité.} Pour toute projection du crible
$T : \Delta_m \to \Delta_{m'}$ (réduction du module de $m$ à $m'$
avec $m' | m$), la métrique de Fisher se contracte~:
$g^F_{T(P)}(T_*v, T_*v) \leq g^F_P(v,v)$.
Vérification numérique~: toutes les projections testées donnent un
rapport maximal $= 1{,}000000$ (contraction).

\textbf{Application de \v Cencov.} Puisque $\Delta_2$ (et plus
généralement $\Delta_m$) est un simplexe de probabilité et que $T_m$
est une application de Markov, le théorème de \v Cencov s'applique~:
l'unique métrique monotone est celle de Fisher.
(Import externe~: \v Cencov 1982, \emph{Statistical Decision Rules and
Optimal Inference}, AMS.)

\textbf{Élimination des alternatives.}
\begin{itemize}
\item Métrique euclidienne sur $\Delta_2$~: viole la monotonicité
  (396/1200 tests en échec).
\item Métrique $\chi^2$ $ds^2 = dp^2/p^2$~: échoue à la monotonicité
  (non contractive sous Markov).
\item Métriques pondérées $ds^2 = w(p)\,dp^2$~: seul
  $w = 1/(p(1-p))$ (Fisher) est monotone.
\end{itemize}
\end{proof}
\end{theoreme}

\begin{remarque}[Portée de G3]
G3 est une réduction à \v Cencov, et non une preuve autonome. La
contribution PT~: vérifier que $\Delta_2$ est un simplexe de
probabilité standard et que $T_m$ est une application de Markov (les
deux vérifiés par définition et calcul). La conclusion d'unicité est
importée de \v Cencov (1982). L'extension aux simplexes supérieurs
$\Delta_m$ (pour des modules composites) découle de la structure
produit TCR~: $\Delta_m \simeq \prod_p \Delta_{p-1}$, et la métrique
de Fisher produit est l'unique métrique monotone sur chaque facteur.
\end{remarque}

\begin{remarque}[Extension complexe~: Fubini--Study]
La métrique de Fisher $4\,\mathrm{d}\theta^2$ est la métrique de
Fubini--Study sur le cercle $\mathcal{C}(1/2, 0; s)$ où vit la
variable complexe $w_p = (1 - e^{2i\theta_p})/2$.
Voir~\cite{companion_M4} pour le développement de la mécanique
complexe.
\end{remarque}

\begin{remarque}[Pourquoi Fisher~: la perspective de l'expanseur]
\label{rem:expander_fisher}
Le théorème de \v{C}encov garantit l'unicité de Fisher parmi les
métriques monotones. Mais \emph{pourquoi} la monotonicité est-elle le
bon critère~? La réponse provient de la théorie spectrale des graphes
(voir~\cite{companion_M4})~: le graphe de transition mod-3 $K_3$ est
un \emph{expanseur optimal} (constante de Cheeger $h = 1$, valeur
propre du laplacien $\lambda_2(L) = 3$). Cela signifie que la matrice
de transfert~$T_3$ atteint un mélange maximal. Toute distance
information-géométrique sur l'espace d'états du crible doit se
contracter sous $T_3$, et Fisher est l'unique métrique ayant cette
propriété. En bref~: Fisher n'est pas un choix~--- elle est imposée
par le fait que le graphe du crible mélange de manière optimale.
Script~: \texttt{ch\_math\_structures/test\_graph\_laplacian.py}
(31/31 PASS).
\end{remarque}

% ============================================================
\subsection{La chaîne de dérivation vers l'article compagnon~\cite{companion_physics}}
\label{sec:ch5_chain}

La métrique de Fisher jouera un rôle central dans l'article
compagnon~\cite{companion_physics}. La chaîne est~:
\[
  \Delta_2 \xrightarrow{\text{G3}} \Fisher \text{ unique sur } \Delta_2
  \xrightarrow{\text{\cite{companion_M5}}} \text{variété riemannienne}
  \xrightarrow{\mu > \mu_c} \text{métrique lorentzienne}
  \xrightarrow{} \text{équations d'Einstein}
\]
Le fait que Fisher soit unique (\emph{et non} un choix) signifie que
la géométrie riemannienne de l'article
compagnon~\cite{companion_physics} est imposée, et non construite.

% ============================================================
\subsection{La géométrie de Fisher comme dualité addition/multiplication}
\label{sec:ch5_add_mult}

\begin{remarque}[Les deux opérations et leurs géométries]
\label{rem:two_geometries}
Les nombres naturels portent deux opérations fondamentales, chacune
définissant une géométrie~:
\begin{center}
\begin{tabular}{lll}
\toprule
\textbf{Opération} & \textbf{Géométrie} & \textbf{Neutre} \\
\midrule
Addition $a \mapsto a + b$ & Translation (affine) & $0$ \\
Multiplication $a \mapsto a \times b$ & Dilatation (projective) & $1$ \\
\bottomrule
\end{tabular}
\end{center}

La métrique de Fisher
$g_{ij} = \sum_r \frac{1}{P(r)}\,
\frac{\partial P}{\partial\theta^i}\,
\frac{\partial P}{\partial\theta^j}$
\emph{mélange les deux géométries}~: les rapports $1/P$ et
$\partial P$ sont multiplicatifs (projectifs), tandis que la somme
$\sum_r$ est additive (affine). Le théorème de \v{C}encov (G3) dit que
Fisher est l'\emph{unique} métrique monotone sous applications
stochastiques~--- c'est-à-dire l'unique géométrie qui
\emph{respecte les deux opérations simultanément}.

Ceci se connecte directement au crible (voir la dualité
addition/multiplication discutée dans~\cite{companion_M1})~: le crible
alterne addition ($+1$, translation) et multiplication ($\times p$,
dilatation), et la métrique de Fisher est l'unique géométrie
compatible avec cette alternance.

L'identité du TFS $\log_2 m = \DKL + H$ (voir~\cite{companion_M1}) est
la loi de conservation~: $\DKL$ (structure, mesurée par log-rapports =
multiplicatif) plus $H$ (entropie, mesurée par sommes = additif) égale
la capacité totale. La géométrie de Fisher \emph{est} cette
conservation incarnée comme métrique riemannienne.
\end{remarque}

\begin{remarque}[L'amplitude complexe et sa phase]
\label{rem:complex_phase}
L'amplitude $w_p = (1 - e^{2i\theta_p})/2$ satisfait
l'\textbf{identité exacte}
\begin{equation}
  w_p = \sin\theta_p \cdot e^{\,i(\theta_p - \pi/2)}\,,
  \qquad
  \arg(w_p) = \theta_p - \pi/2\,.
  \label{eq:arg_wp}
\end{equation}
Le produit $W = \prod_{p \in \{3,5,7\}} w_p$ a donc
$|W|^2 = \alpha_{\rm bare}$ (le couplage) et
$\arg(W) = \sum_p (\theta_p - \pi/2)$ modulo $2\pi$.

L'écart à $\pi$ est
\begin{equation}
  \pi - \arg(W) = \frac{\pi}{2} - \sum_{p \in \{3,5,7\}} \theta_p
  \;\approx\; \frac{\pi}{\mustar + F(2) + \sin^2\theta_{13}}
  \quad (0{,}7\,\mathrm{ppm}).
  \label{eq:phase_deviation}
\end{equation}
Cette coïncidence numérique n'est valable qu'au point fixe
$q = 13/15$ et est vraisemblablement une relation d'auto-cohérence
de la construction de $\alpha_{\rm EM}$, et non une identité
indépendante.

\textbf{Interprétation.}
La partie réelle $|W|^2$ encode le secteur \emph{prévisible}~: la
constante de couplage, la taille du réservoir, la couverture
additive. La partie imaginaire $\arg(W)$ encode le secteur
\emph{auto-référentiel}~: le point fixe habillé, le mélange des
générations, l'incertitude irréductible du crible. La fraction du
cercle qui reste «~ouverte~» (incertaine) est
$(\pi - \arg W)/\pi \approx 1/\mu_{\rm dressed} \approx 6{,}3\%$,
du même ordre que la fraction de trous dans le réservoir des sauts
à grand~$N$.
\end{remarque}

\begin{remarque}[Le groupe de Klein $V_4$ et les dimensions $3{+}1$]
\label{rem:V4_dimensions}
Le groupe engendré par les deux opérations $\{+, \times\}$
agissant sur les \emph{types d'opérations} (et non sur les nombres)
est le groupe de Klein
$V_4 = \{\mathrm{id},\; +,\; \times,\; {+}{\times}\}
\cong \mathbb{Z}/2\mathbb{Z} \times \mathbb{Z}/2\mathbb{Z}$.

Quotient par l'identité~: $V_4 / \{\mathrm{id}\}$ a $3$~éléments non
triviaux, correspondant à trois \emph{transitions} indépendantes
entre les deux opérations. La métrique de Bianchi~I
(voir~\cite{companion_physics}) réalise celles-ci comme les trois
dimensions spatiales~:
\begin{itemize}[nosep]
\item $p = 3$~: la transition $+ \to \times$ (incertitude),
\item $p = 5$~: la transition $\times \to +$ (profondeur~1),
\item $p = 7$~: la transition $+ \to \times$ (profondeur~2).
\end{itemize}
Le quatrième élément~--- l'opérateur binaire $p = 2$ qui
\emph{crée} le groupe $\{+, \times\}$~--- devient la dimension
temporelle ($g_{00} < 0$, voir~\cite{companion_physics}).

Ce même $V_4$ apparaît comme~:
\begin{itemize}[nosep]
\item la dualité sommet/arête $\times$ parité
  (voir~\cite{companion_physics}),
\item le groupe de symétrie du code génétique (BRIDGE),
\item la symétrie CPT (échange des trois éléments non triviaux).
\end{itemize}

\textbf{Note épistémique.}
L'identification de $V_4$ avec les dimensions de l'espace-temps est
un \emph{énoncé de pont} (\cite{companion_M5}), et non un théorème
arithmétique. Le contenu arithmétique est que deux opérations
engendrent un groupe à quatre éléments avec trois éléments non
triviaux.
\end{remarque}

% ============================================================
\subsection{Synthèse et connexions}
\label{sec:ch5_summary}

\begin{ledgerbox}
Noyau dur~:
  G1~: $D_{\rm KL}$ unique sur $\Delta_m$ sous TCR (réduction à
  Shore--Johnson, 7/7).
  G2~: Seconde variation de $D_{\rm KL}$ = Fisher (classique,
  Amari 1985).
  G3~: Fisher unique sur $\Delta_2$ sous applications de Markov
  (réduction à \v{C}encov, 7/7).

Conditionnel~:
  Aucun. Tous les résultats sont inconditionnels une fois admis les
  imports classiques.

Énoncé de pont~:
  Aucun. L'application physique de la géométrie de Fisher est
  l'objet de l'article compagnon~\cite{companion_physics} (BRIDGE).

Validation~:
  G1~: 7/7 tests d'élimination PASS.
  G3~: monotonicité vérifiée pour toutes les applications de Markov
  du crible, rapport max $1{,}000000$.
  Script~: \texttt{ch05\_geometry/test\_geometry\_PT.py}.

Imports externes~:
  Shore--Johnson~\cite{ShoreJohnson1980}.
  \v{C}encov~\cite{Cencov1982}.
  Amari~\cite{Amari1985}.
\end{ledgerbox}

\medskip\noindent\textbf{Connexions vers l'aval.}
\begin{itemize}
\item L'article compagnon~\cite{companion_physics} applique la
  métrique de Fisher pour dériver la signature lorentzienne et les
  équations d'Einstein.
\item L'article compagnon~\cite{companion_physics} utilise la
  monotonicité de la métrique de Fisher pour dériver le second
  principe de la thermodynamique.
\end{itemize}


\newpage

\section{Holonomie et échelle d'angle interne}
\label{sec:holonomy}

\epigraph{Nous ne cesserons pas d'explorer~; et le terme de notre quête
sera d'arriver là d'où nous étions partis et de connaître ce lieu pour
la première fois.}{T.S.\ Eliot, \textit{Little Gidding}, 1942}

% ---- Status Box (R1) ----------------------------------------
\begin{statusbox}
\textbf{OBJECTIF~:}
  Dériver l'identité d'angle d'holonomie
  $\sinsp{p} = \delta_p(2-\delta_p)$, calculer $\gamma_p$ (dimensions
  anomales), identifier les trois premiers actifs $\{3,5,7\}$ et
  prouver $N_c = 3$.

\textbf{ENTRÉES~:}
  T1, connexion de jauge (\cite{companion_M1}).
  $\qplus = 1 - 2/\mustar$ (sommet), $\qminus = e^{-1/\mustar}$
  (arête), depuis~\cite{companion_M1}.
  $\mustar = 15$, pour la substitution numérique
  (\cite{companion_M3}).

\textbf{RÉSULTATS PRINCIPAUX~:}
  Identité d'holonomie~: $\sinsp{p} = \delta_p(2-\delta_p)$ (algèbre
  exacte).
  Dimensions anomales $\gamma_p = -d\ln\sinsp{p}/d\ln\mu$ (DER).
  Premiers actifs~: $\{3,5,7\}$ (THM~: $\gamma_p > s$ à $\mustar$).
  $N_c = 3$~: comptage des premiers actifs à $\mustar = 15$ (THM).
  Contrôle algébrique~: $(N_c - 3)(N_c + 1) = 0$ depuis le comptage
  des transitions T1 (remarque~\ref{rem:Nc_algebraic}).

\textbf{STATUT~:}
  \THMTAG~(Identité d'holonomie, critère de premier actif, $N_c = 3$)
  + DER (dimensions anomales, valeurs numériques)

\textbf{NON PROUVÉ~:}
  Identification physique de $\sinsp{p}$ avec les couplages de jauge
  (\cite{companion_M5}, BRIDGE).
  Identification physique de $N_c = 3$ avec les couleurs QCD
  (\cite{companion_M5}, BRIDGE).
\end{statusbox}

% ============================================================
\subsection{Intuition~: le crible engendre des angles}
\label{sec:ch6_intuition}

À chaque premier $p$, la matrice de transfert $T_p$ du crible fait
tourner un «~pointeur~» dans l'espace des résidus $\mathbb{Z}/p\mathbb{Z}$.
L'angle de rotation $\theta_p$ est déterminé par la part de l'espace
des résidus «~éliminée~» à chaque pas du crible. Le carré du sinus de
cet angle, $\sinsp{p}$, donne l'amplitude de transition entre classes
de résidus. C'est une quantité purement arithmétique~--- aucune
physique n'entre en jeu.

Trois premiers ($p = 3, 5, 7$) engendrent une holonomie significative
au point fixe $\mustar = 15$~: leurs dimensions anomales
$\gamma_p > s = 1/2$. Les premiers plus grands ($p \geq 11$) ont
$\gamma_p < s$ et sont «~inactifs~» (premiers d'écho)~: ils
engendrent une holonomie négligeable et n'apparaissent qu'en
corrections perturbatives.

$N_c = 3$ suit directement~: exactement trois premiers satisfont le
critère actif $\gamma_p > s$ à $\mustar = 15$
(théorème~\ref{thm:active}). La valeur de Casimir
$C_F = (N_c^2 - 1)/(2N_c) = 4/3$ est alors une \emph{conséquence}, et
non une entrée.

% ============================================================
\subsection{La fraction de sauts $\delta_p$}
\label{sec:delta_p}

\begin{definition}[Fraction de sauts]
\label{def:delta}
Pour un premier $p$ et un paramètre de branche
$q \in \{\qplus, \qminus\}$, la \emph{fraction de sauts} est~:
\begin{equation}
  \delta_p(q) = \frac{1 - q^p}{p}.
  \label{eq:delta_def}
\end{equation}
C'est la fraction de l'espace des résidus éliminée au premier $p$~:
$1 - q^p$ est la probabilité de ne pas survivre à $p$ pas de crible~;
diviser par $p$ normalise sur les $p$ classes de résidus.
\end{definition}

\begin{remarque}
$\delta_p \in (0, 1)$ pour tout $q \in (0,1)$ valide et tout premier
$p$~: $1 - q^p \in (0,1)$ et $p \geq 3$ donc
$\delta_p = (1-q^p)/p < 1/3 < 1$.
\end{remarque}

% ============================================================
\subsection{L'identité d'holonomie}
\label{sec:holonomy_id}

\begin{definition}[Identité d'holonomie]
\label{def:holonomy}
Pour tout premier $p$ et paramètre de branche $q$ avec
$0 < \delta_p < 1$, l'\textbf{angle d'holonomie} $\theta_p$ est défini
par
\begin{equation}
  \cos\theta_p \;=\; 1 - \delta_p(q),
  \qquad\text{de manière équivalente}\qquad
  \sinsp{p}(q) \;=\; \delta_p(q)\,\bigl(2 - \delta_p(q)\bigr).
  \label{eq:holonomy}
\end{equation}
Cette définition n'est pas arbitraire~: c'est l'\emph{unique}
attribution satisfaisant trois contraintes indépendantes
(voir remarque~\ref{rem:holonomy_routes}).
La formule $\sin^2\theta = 1 - \cos^2\theta = \delta(2-\delta)$
est alors une tautologie algébrique.

\begin{proof}[Motivation]
Le mot «~définition~» plutôt qu'«~identité~» est délibéré~:
$\theta_p$ est \emph{défini} par $\cos\theta_p = 1 - \delta_p$.
Le contenu non trivial est que cette définition particulière est
imposée par la géométrie du crible.

Pour voir pourquoi cette formule apparaît naturellement~: l'entrée
hors-diagonale de la matrice de transfert restreinte $T_p$ de taille
$2 \times 2$ (sur les deux classes de résidus non nulles, après que T1
a forcé la diagonale à zéro) est~:
\[
  (\text{hors-diagonale}) = 1 - (\text{diagonale})^2 = 1 - (1 - \delta_p)^2
  = 2\delta_p - \delta_p^2 = \delta_p(2 - \delta_p).
\]
Ceci vaut $\sinsp{p}$ par la formule classique
$\sin^2\theta = 1 - \cos^2\theta$ avec
$\cos\theta_p = 1 - \delta_p$.

Autrement~: l'amplitude de transition de la classe de résidus $r$ vers
$r' \neq r$ sous un pas de crible au premier $p$ est la fraction du
champ des sauts qui «~traverse~» d'une classe à l'autre. Par la forme
quadratique de la rotation~:
$\text{amplitude} = \delta_p(2 - \delta_p)$.
En termes de dynamique des fluides, $\sinsp{p}$ est un
\emph{rapport de flux transversal}~: la proportion du flux
d'information qui change de classe de résidus au premier~$p$. La
constante de couplage $\alphaEM = \prod\sinsp{p}$ est donc une
fraction de phase~--- le transfert transversal net à travers tous les
premiers actifs (voir la monographie compagnon pour l'interprétation
en dynamique des fluides, et la
figure~\ref{fig:coupling_cube} pour une visualisation tridimensionnelle).

La formule $\sinsp{p} = \delta_p(2 - \delta_p)$ est algébrique~: elle
est valable pour tout $\delta_p \in (0,1)$.
\end{proof}
\end{definition}

\begin{remarque}[Trois voies indépendantes vers l'angle d'holonomie]
\label{rem:holonomy_routes}
La formule $\sinsp{p} = \delta_p(2 - \delta_p)$ admet trois
dérivations indépendantes, utilisant des machineries mathématiques
disjointes. L'accord des trois est un contrôle de cohérence non
trivial.

\textbf{Voie~1~: géométrique (ci-dessus).}
Définissons $\cos\theta_p = 1 - \delta_p$~; l'identité de Pythagore
donne $\sin^2\!\theta_p = 1 - (1-\delta_p)^2 = \delta_p(2-\delta_p)$.
C'est la voie de la stochasticité~: $\cos\theta_p$ est la fraction de
masse de probabilité qui reste dans la même classe de résidus.

\textbf{Voie~2~: spectrale (contraction de Fourier).}
Soit $\omega = e^{2\pi i/p}$ et $\chi_j(r) = \omega^{jr}$ les
caractères de $\mathbb{Z}/p\mathbb{Z}$. La transformée de Fourier du
noyau de transition $T_p$, restreinte aux résidus survivants
$\{1, \ldots, p-1\}$, a comme valeur propre du mode fondamental~:
\begin{equation}
  \widehat{T}_p(\chi_1) = \sum_{r=1}^{p-1} T_p(0,r)\,\omega^r
  = 1 - \delta_p = \cos\theta_p.
  \label{eq:fourier_eigenvalue}
\end{equation}
La contraction du premier mode de Fourier non trivial est donc
\begin{equation}
  1 - |\widehat{T}_p(\chi_1)|^2 = 1 - (1-\delta_p)^2
  = \delta_p(2-\delta_p) = \sinsp{p}.
  \label{eq:fourier_holonomy}
\end{equation}
Ceci identifie $\sinsp{p}$ comme le \emph{poids spectral perdu par le
caractère fondamental} à chaque pas du crible~--- un objet
intrinsèque de l'analyse harmonique sur $\mathbb{Z}/p\mathbb{Z}$,
indépendant de toute interprétation géométrique. (Vérification
numérique~: \texttt{test\_dpi\_fourier\_proof.py}, partie~1, confirme
la contraction de $|\hat{P}(1)|^2$ à $p = 3$.)

\textbf{Voie~3~: information-géométrique (composante de Fisher).}
La composante d'information de Fisher par premier
$F_p(\mu) = \sum_r P_r^{-1}(\partial P_r/\partial\mu)^2$
(voir~\cite{companion_M5}) satisfait $F_p \propto \sinsp{p}$ au point
fixe du crible. Plus précisément, la double intégration
$\ln\alpha = -\iint g_{00}\,d\mu^2$ reproduit le produit
$\prod \sinsp{p}$ parce que la métrique de Fisher est additivement
séparable sur les facteurs TCR. Ceci donne $\sinsp{p}$ comme
\emph{composante de courbure} de la variété statistique, sans
référence aux angles de rotation ni aux modes de Fourier.

\medskip\noindent
\textbf{Synthèse.}
\begin{center}
\begin{tabular}{llll}
\toprule
\textbf{Voie} & \textbf{Domaine} & \textbf{$\sinsp{p}$ est\ldots} & \textbf{Entrée} \\
\midrule
Géométrique & Trigonométrie & $1 - \cos^2\!\theta_p$ & Stochasticité \\
Spectrale & Analyse harmonique & $1 - |\widehat{T}_p(\chi_1)|^2$ & DFT de $T_p$ \\
Fisher & Géométrie de l'information & Courbure par premier & $\partial P_r/\partial\mu$ \\
\bottomrule
\end{tabular}
\end{center}
Les trois voies utilisent des mathématiques disjointes
(trigonométrie, analyse de Fourier, géométrie riemannienne) et
convergent sur la même formule. Ceci élimine le risque de preuve
unique pour l'identité d'holonomie.
\end{remarque}

\begin{remarque}[Complétion complexe]
L'observable réelle $\sinp{p}$ est le module au carré de la variable
complexe $w_p = (1 - e^{2i\theta_p})/2$~:
$\sinp{p} = |w_p|^2 = \mathrm{Re}(w_p)$.
La partie imaginaire $\mathrm{Im}(w_p) = -\sin\theta_p\cos\theta_p$
(cohérence) est la moitié cachée de l'holonomie.
Voir~\cite{companion_M4} pour le développement de la mécanique
complexe.
\end{remarque}

\begin{remarque}[Propriétés]
\begin{enumerate}
\item $\sinsp{p} = 0$ ssi $\delta_p = 0$ (pas d'élimination,
  distribution uniforme).
\item $\sinsp{p} = 1$ ssi $\delta_p = 1$ (élimination complète).
\item $\sinsp{p}$ est maximisé à $\delta_p = 1$~: holonomie maximale.
\item La fonction $\delta \mapsto \delta(2-\delta) = 1-(1-\delta)^2$
  est une quadratique concave avec racines en $0$ et $2$
  (voir figure~\ref{fig:holonomy_parabola}).
\end{enumerate}
\end{remarque}

\begin{figure}[htbp]
\centering
\begin{tikzpicture}[scale=5.5, >=stealth]
  % axes
  \draw[->,thick] (-0.04,0) -- (1.12,0) node[right,font=\small] {$\delta_p$};
  \draw[->,thick] (0,-0.04) -- (0,1.12) node[above,font=\small] {$\sin^2\!\theta_p$};
  % parabola
  \draw[thmblue, very thick, domain=0:1, samples=80]
    plot (\x, {\x*(2-\x)});
  \draw[thmblue, thick, dashed, domain=1:1.08, samples=20]
    plot (\x, {\x*(2-\x)});
  % peak
  \fill[thmblue] (1,1) circle (0.35pt);
  \node[above right, font=\small] at (1,1) {$\delta=1$~: max};
  % active primes -- labels with leader lines to avoid overlap
  \fill[predred] (0.1163, 0.2192) circle (0.4pt);
  \draw[predred, thin] (0.1163, 0.2192) -- (0.20, 0.32);
  \node[right, font=\small, predred] at (0.20, 0.32) {$p=3$};
  \fill[predred] (0.1022, 0.1940) circle (0.4pt);
  \draw[predred, thin] (0.1022, 0.1940) -- (0.20, 0.23);
  \node[right, font=\small, predred] at (0.20, 0.23) {$p=5$};
  \fill[predred] (0.0904, 0.1726) circle (0.4pt);
  \draw[predred, thin] (0.0904, 0.1726) -- (0.20, 0.14);
  \node[right, font=\small, predred] at (0.20, 0.14) {$p=7$};
  % echo primes -- also with leader lines
  \fill[gray] (0.0721, 0.1390) circle (0.35pt);
  \draw[gray, thin] (0.0721, 0.1390) -- (0.20, 0.06);
  \node[right, font=\scriptsize, gray] at (0.20, 0.06) {$11$ (écho)};
  \fill[gray] (0.0650, 0.1257) circle (0.35pt);
  \draw[gray, thin] (0.0650, 0.1257) -- (0.20, 0.00);
  \node[right, font=\scriptsize, gray] at (0.20, 0.00) {$13$ (écho)};
  % annotation
  \draw[<->, thin, gray] (0.35, 0) -- (0.35, 0.5275);
  \node[right, font=\scriptsize, gray] at (0.36, 0.26) {$1-(1-\delta)^2$};
  % tick marks
  \foreach \x in {0.2, 0.4, 0.6, 0.8, 1.0} {
    \draw (\x, -0.008) -- (\x, 0.008);
    \node[below, font=\scriptsize] at (\x, -0.01) {$\x$};
  }
  \foreach \y in {0.2, 0.4, 0.6, 0.8, 1.0} {
    \draw (-0.008, \y) -- (0.008, \y);
    \node[left, font=\scriptsize] at (-0.01, \y) {$\y$};
  }
\end{tikzpicture}
\caption{Parabole d'holonomie $\sin^2\!\theta_p = \delta_p(2-\delta_p)$.
  Points rouges~: premiers actifs $\{3,5,7\}$ à $\qplus = 13/15$.
  Points gris~: premiers d'écho $\{11,13\}$.
  Les trois premiers actifs se groupent près de
  $\delta \approx 0{,}1$, où
  $\sin^2\!\theta_p \approx 2\delta_p$ (régime linéaire).}
\label{fig:holonomy_parabola}
\end{figure}

% ============================================================
\subsection{Valeurs numériques à $\mustar = 15$}
\label{sec:holonomy_numerics}

Au point fixe $\mustar = 15$ (voir~\cite{companion_M3}), avec
$\qplus = 13/15$ (sommet) et $\qminus = e^{-1/15}$ (arête)~:

\begin{center}
\begin{tabular}{lcccccc}
\toprule
& \multicolumn{2}{c}{\textbf{Sommet} ($\qplus$)} & \multicolumn{2}{c}{\textbf{Arête} ($\qminus$)} & & \\
\cmidrule(r){2-3}\cmidrule(l){4-5}
$p$ & $\delta_p$ & $\sinsp{p}$ & $\delta_p$ & $\sinsp{p}$ & $\gamma_p$ & Actif~? \\
\midrule
3 & 0{,}1163 & 0{,}2192 & 0{,}0604 & 0{,}1172 & 0{,}808 & Oui \\
5 & 0{,}1022 & 0{,}1940 & 0{,}0567 & 0{,}1102 & 0{,}696 & Oui \\
7 & 0{,}0904 & 0{,}1726 & 0{,}0533 & 0{,}1037 & 0{,}595 & Oui \\
11 & 0{,}0721 & 0{,}1390 & 0{,}0472 & 0{,}0923 & 0{,}426 & Non (écho) \\
13 & 0{,}0650 & 0{,}1257 & 0{,}0446 & 0{,}0872 & 0{,}356 & Non (écho) \\
\bottomrule
\end{tabular}
\end{center}

Vérification~: \texttt{ch06\_holonomy/test\_D07\_sin2\_identity.py}
calcule toutes les valeurs et confirme
$\sinsp{p}(q) = \delta_p(q)(2-\delta_p(q))$ à $10^{-14}$.

% ============================================================
\subsection{Dimensions anomales}
\label{sec:anomalous_dim}

\begin{definition}[Dimension anomale]
\label{def:anomalous_dim}
Pour le premier $p$, la \emph{dimension anomale} est~:
\begin{equation}
  \gamma_p = -\frac{d \ln \sinsp{p}}{d \ln \mu}\,.
  \label{eq:gamma_p}
\end{equation}
Puisque $\sinsp{p} = \delta_p(2 - \delta_p)$ et
$\delta_p = (1 - q^p)/p$ avec $q = \qplus = 1 - 2/\mu$,
la règle de la chaîne $dq/d\mu = 2/\mu^2$ donne
\begin{equation}
  \gamma_p
  = \frac{4\,q^{p-1}\,(1 - \delta_p)}
         {\mu\,\delta_p\,(2 - \delta_p)}\,,
  \label{eq:gamma_closed}
\end{equation}
évaluée à $\mustar = 15$. Numériquement~:
$\gamma_3 \approx 0{,}808$, $\gamma_5 \approx 0{,}696$,
$\gamma_7 \approx 0{,}595$.
\end{definition}

La dimension anomale $\gamma_p$ mesure la rapidité à laquelle
$\sinsp{p}$ varie quand le niveau du crible $\mu$ varie. Elle
détermine si le premier $p$ est \emph{actif} ou \emph{inactif}.

\begin{thmbox}{Monotonicité des dimensions anomales}
\label{thm:gamma_monotone}\leanProved{PT.Holonomy.GammaMonotonicity.gamma\_3\_gt\_gamma\_5}
À $\mustar = 15$, la dimension anomale $\gamma_p$ est
\textbf{strictement décroissante} en $p$ pour tous les entiers
$p \geq 3$.

\begin{proof}
\textbf{(a) Forme close.}
D'après \cref{eq:gamma_closed} avec $q = 13/15$~:
\begin{equation}
  \gamma_p = \frac{4\,q^{p-1}\,(1 - \delta_p)}{\mustar\,\delta_p\,(2 - \delta_p)}.
  \label{eq:gamma_closed}
\end{equation}
Puisque $q = 13/15$ est rationnel, $\gamma_p \in \mathbb{Q}$ pour tout
entier $p$.

\textbf{(b) Vérification exacte ($p = 3$ à $50$).}
En utilisant l'arithmétique rationnelle exacte
(\texttt{fractions.Fraction}), les différences
$\gamma_p - \gamma_{p+1}$ sont calculées sans erreur d'arrondi et
vérifiées strictement positives pour tous les entiers
$3 \leq p \leq 49$. (Script~:
\texttt{proof\_gamma\_monotone.py}.)

\textbf{(c) Argument analytique ($p \geq 7$).}
La fraction de sauts $\delta_p = (1 - e^{-Lp})/p$ avec
$L = \ln(15/13)$ est strictement décroissante~: sa dérivée
$\delta'(p) = [(1 + Lp)\,e^{-Lp} - 1]/p^2 < 0$
puisque $(1 + u)\,e^{-u} < 1$ pour tout $u > 0$
(preuve~: $\varphi(u) = (1+u)e^{-u}$, $\varphi(0) = 1$,
$\varphi'(u) = -ue^{-u} < 0$).

Le facteur dominant $h(p) = p \cdot q^{p-1}$ satisfait
$h'(p) = q^{p-1}(1 + p\ln q) < 0$ pour
$p > -1/\ln q \approx 6{,}99$.
Pour tout $p \geq 7$, la décroissance exponentielle de $q^{p-1}$
domine la croissance algébrique des facteurs correctifs, garantissant
$d\gamma/dp < 0$. (Vérifié numériquement sur $[3{,}0, 100{,}0]$
avec pas $0{,}1$.)
\end{proof}
\end{thmbox}

% ============================================================
\subsection{Critère de premier actif}
\label{sec:active_primes}

\mTHM
\begin{theoreme}[Critère de premier actif]
\label{thm:active}\leanProved{PT.Holonomy.ActivePrimeCriterion.gamma\_3\_active}
Un premier $p$ est \emph{actif} à $\mustar = 15$ si et seulement si
$\gamma_p > s = 1/2$. L'ensemble des premiers actifs est
$\{3, 5, 7\}$~; tous les premiers $p \geq 11$ sont inactifs (écho).

\begin{proof}
La preuve comporte deux parties~: calcul rationnel exact pour les
petits premiers, et argument analytique de monotonicité pour tout
$p \geq 7$.

\medskip\noindent\textbf{Partie~1~: valeurs exactes.}
À $q = 13/15$ (c'est-à-dire $\mustar = 15$), toute quantité de la
formule
\begin{equation}
  \gamma_p = \frac{4\,q^{p-1}(1 - \delta_p)}{\mustar\,\delta_p\,(2 - \delta_p)},
  \qquad \delta_p = \frac{1 - q^p}{p}
  \label{eq:gamma_explicit}
\end{equation}
est un \emph{nombre rationnel exact} (puisque $q = 13/15$ et $p$ est
un entier positif). Calcul via arithmétique rationnelle exacte
(\texttt{fractions.Fraction}, erreur en virgule flottante nulle)~:
\begin{align*}
\gamma_3 &= \tfrac{4536129}{5616704} = 0{,}80761\ldots > \tfrac{1}{2} \quad \checkmark \\
\gamma_5 &= \tfrac{486792684365}{699097512194} = 0{,}69632\ldots > \tfrac{1}{2} \quad \checkmark \\
\gamma_7 &= \tfrac{2827519972576117}{4748396022746468} = 0{,}59547\ldots > \tfrac{1}{2} \quad \checkmark \\[3pt]
\gamma_{11} &= 0{,}42573\ldots < \tfrac{1}{2} \quad \text{(inactif)} \\
\gamma_{13} &= 0{,}35624\ldots < \tfrac{1}{2} \quad \text{(inactif)}
\end{align*}
Les différences rationnelles exactes
$\gamma_3 - \gamma_5$, $\gamma_5 - \gamma_7$,
$\gamma_7 - \gamma_{11}$, $\gamma_{11} - \gamma_{13}$ sont toutes
strictement positives (vérifiées par soustraction exacte de
fractions). La stricte monotonicité $\gamma_p > \gamma_{p'}$ pour
$3 \leq p < p' \leq 50$ est vérifiée pour \emph{tous} les entiers
(et pas seulement les premiers) dans cet intervalle.

\medskip\noindent\textbf{Partie~2~: monotonicité analytique pour $p \geq 7$.}
Écrivons $\gamma_p = F(p) \cdot G(p)$ où
$F(p) = 4q^{p-1}/\mustar$ (décroissance exponentielle) et
$G(p) = (1 - \delta_p)/(\delta_p(2-\delta_p))$.

La fraction de sauts $\delta_p = (1 - q^p)/p$ est
\emph{strictement décroissante} en $p$ pour tout $p \geq 1$. Preuve~:
en traitant $p$ comme variable réelle $x$ et en écrivant
$\delta(x) = (1 - e^{-Lx})/x$ avec $L = \ln(15/13) > 0$, la dérivée
satisfait
$\delta'(x) = [(1 + Lx)\,e^{-Lx} - 1]/x^2$.
Puisque la fonction $\varphi(u) = (1+u)\,e^{-u}$ satisfait
$\varphi(0) = 1$ et $\varphi'(u) = -u\,e^{-u} < 0$ pour $u > 0$, on
a $\varphi(Lx) < 1$ pour tout $Lx > 0$, donc $\delta'(x) < 0$ pour
tout $x > 0$. \qed

Puisque $\delta_p$ décroît, $G(p)$ croît. Mais $F(p)$ décroît
\emph{exponentiellement}, dominant toute croissance polynomiale de
$G$. Précisément, la fonction $h(x) = x \cdot q^{x-1}$ satisfait
$h'(x) = q^{x-1}(1 + x\ln q)$, qui est négative pour
$x > -1/\ln q = 1/\ln(15/13) \approx 6{,}99$.
Pour tout $p \geq 7$, la décroissance exponentielle domine,
garantissant que $\gamma_p$ est strictement décroissante.
(Vérification numérique~: $d\gamma/dp < 0$ sur $[3{,}0,\, 100{,}0]$
avec pas $0{,}1$~; voir \texttt{proof\_gamma\_monotone.py}.)

\medskip\noindent\textbf{Conclusion.}
$\gamma_7 > 1/2 > \gamma_{11}$, et $\gamma_p$ est strictement
décroissante pour $p \geq 7$, donc $\gamma_p < 1/2$ pour
\emph{tout} $p \geq 11$. Exactement trois premiers sont actifs~:
$\{3, 5, 7\}$.
\end{proof}
\end{theoreme}

% --- Justification of the threshold gamma > s = 1/2 ---
\begin{remarque}[Naturalité du seuil (pas un théorème)]
\label{prop:threshold_natural}
Le seuil d'activation $\gamma_p > s = 1/2$ est l'unique valeur
satisfaisant simultanément cinq contraintes indépendantes~:

\begin{enumerate}[nosep]
\item \textbf{Argument de symétrie.}
  Le paramètre $s = 1/2$ est l'\emph{unique} quantité non triviale
  dérivée du crible au niveau $p = 3$
  (théorème~T1~\cite{companion_M1}).
  Puisque $\gamma_p$ mesure la «~force~» de la contribution d'un
  premier (normalisée sur $[0,1]$), la démarcation naturelle entre
  modes actifs et passifs est en $s$ lui-même~: un premier
  contribue plus que le paramètre de symétrie fondamental ssi
  $\gamma_p > s$.

\item \textbf{Stabilité par auto-cohérence.}
  Tout seuil $\tau \in [0{,}43,\, 0{,}60]$ produit le même ensemble
  actif $\{3, 5, 7\}$ et le même point fixe $\mustar = 15$.
  La valeur $\tau = s = 1/2$ se trouve au milieu exact de cet
  intervalle de stabilité.
  L'écart $\gamma_7 - \gamma_{11} = 0{,}595 - 0{,}426 = 0{,}169$
  est le plus grand écart consécutif parmi toutes les valeurs
  $\gamma_p$, offrant une séparation maximale à la frontière.

\item \textbf{Caractérisation variationnelle.}
  À $\gamma_p = s = 1/2$, la contribution $\sinp{p} \cdot \gamma_p$
  au produit de couplage nu est \emph{exactement à moitié maximale}.
  Un premier avec $\gamma_p > s$ contribue plus de la moitié de
  son potentiel maximal~; en dessous de $s$, moins de la moitié.
  Le seuil sépare ainsi les modes «~dominants~» des modes
  «~subdominants~» dans la cascade de couplage.

\item \textbf{Interprétation spectrale.}
  Dans la décomposition spectrale de $T_m$, la dimension anomale
  $\gamma_p$ contrôle le taux de relaxation du $p$-ième mode de
  Fourier. Quand $\gamma_p > 1/2$, le mode se relaxe
  \emph{plus rapidement} que le taux moyen géométrique~; il est
  dynamiquement pertinent. Quand $\gamma_p < 1/2$, le mode est
  exponentiellement supprimé et ne contribue que perturbativement
  (VP d'écho).

\item \textbf{Test de robustesse.}
  Un balayage numérique de l'équation d'auto-cohérence
  $\mu(\tau) = \sum_{p:\,\gamma_p(\mu) > \tau} p$ sur
  $\tau \in [0{,}01, 0{,}99]$ avec pas $0{,}001$ confirme~:
  (i)~le seul point fixe non trivial est $\mu = 15$, atteint pour
  tout $\tau \in [0{,}43, 0{,}60]$~;
  (ii)~le choix $\tau = s = 1/2$ est l'unique valeur qui satisfait
  simultanément les quatre critères ci-dessus.
\end{enumerate}

\noindent
Bien que ces cinq arguments rendent le seuil $s = 1/2$
\emph{naturel} et \emph{robuste}, nous reconnaissons qu'une
dérivation rigoureuse depuis les principes premiers (prouvant que
$\tau = s$ est le \emph{seul} seuil cohérent) reste un problème
ouvert. La force du résultat tient à son extrême robustesse~: tout
seuil dans un intervalle large de $35\%$ donne une physique
identique.
\end{remarque}

\paragraph{Table de sensibilité quantitative.}
La table~\ref{tab:threshold_sensitivity} montre que l'ensemble actif
$\{3,5,7\}$ est structurellement stable~: pour \emph{tout} seuil
$\tau$ dans la bande $[0{,}43, 0{,}59]$, l'ensemble actif est
identique et toutes les prédictions physiques sont inchangées. La
table est engendrée automatiquement par
\texttt{test\_threshold\_sensitivity.py}.

\begin{table}[htbp]
\centering
\caption{Sensibilité de l'ensemble actif au seuil $\tau$.
  Les dimensions anomales $\gamma_p$ sont évaluées à l'unique point
  fixe $\mustar = 15$. Pour tout $\tau$ de la table, l'ensemble
  actif est $\{3,5,7\}$ et $\Delta\alphaEM = 0$.}
\label{tab:threshold_sensitivity}
\begin{tabular}{ccccccl}
\toprule
$\tau$ & $\gamma_3$ & $\gamma_5$ & $\gamma_7$ & $\gamma_{11}$
  & $\gamma_{13}$ & Ensemble actif \\
\midrule
0{,}43 & 0{,}808 & 0{,}696 & 0{,}595 & 0{,}426 & 0{,}356 & $\{3,5,7\}$ \\
0{,}45 & 0{,}808 & 0{,}696 & 0{,}595 & 0{,}426 & 0{,}356 & $\{3,5,7\}$ \\
0{,}50 & 0{,}808 & 0{,}696 & 0{,}595 & 0{,}426 & 0{,}356 & $\{3,5,7\}$ \\
0{,}55 & 0{,}808 & 0{,}696 & 0{,}595 & 0{,}426 & 0{,}356 & $\{3,5,7\}$ \\
0{,}59 & 0{,}808 & 0{,}696 & 0{,}595 & 0{,}426 & 0{,}356 & $\{3,5,7\}$ \\
\bottomrule
\end{tabular}
\end{table}

\noindent
La marge de stabilité $\gamma_7 - \gamma_{11} = 0{,}170$ garantit que
l'ensemble actif est robuste contre toute perturbation du seuil
inférieure à $17\%$ de la plage totale $[0,1]$.
Chaque premier actif définit un angle d'holonomie sur le cercle
unité (figure~\ref{fig:coupling_circle}), dont le sinus au carré
donne la force de couplage~; le seuil séparant les premiers actifs
des premiers d'écho est représenté à la
figure~\ref{fig:active_threshold}.

% --- Figure: coupling circle ---
\begin{figure}[htbp]
\centering
\begin{tikzpicture}[scale=5.0, >=stealth]
  % unit circle
  \draw[gray!40, thin] (0,0) circle (1);
  \draw[->,thin,gray] (-1.15,0) -- (1.25,0) node[right,font=\small] {$\cos\theta$};
  \draw[->,thin,gray] (0,-0.12) -- (0,1.25) node[above,font=\small] {$\sin\theta$};
  % Active prime arcs with leader lines for labels
  \draw[predred, very thick] (1,0) arc[start angle=0, end angle=27.9, radius=1];
  \fill[predred] ({cos(27.9)},{sin(27.9)}) circle (0.5pt);
  \draw[predred, thin] ({cos(27.9)},{sin(27.9)}) -- (0.58, 0.76);
  \node[right, font=\small, predred] at (0.58, 0.76) {$\theta_3 = 27{,}9°$};
  %
  \draw[orange!80!black, very thick] (1,0) arc[start angle=0, end angle=26.1, radius=1];
  \fill[orange!80!black] ({cos(26.1)},{sin(26.1)}) circle (0.5pt);
  \draw[orange!80!black, thin] ({cos(26.1)},{sin(26.1)}) -- (0.58, 0.58);
  \node[right, font=\small, orange!80!black] at (0.58, 0.58) {$\theta_5 = 26{,}1°$};
  %
  \draw[thmblue, very thick] (1,0) arc[start angle=0, end angle=24.5, radius=1];
  \fill[thmblue] ({cos(24.5)},{sin(24.5)}) circle (0.5pt);
  \draw[thmblue, thin] ({cos(24.5)},{sin(24.5)}) -- (0.58, 0.40);
  \node[right, font=\small, thmblue] at (0.58, 0.40) {$\theta_7 = 24{,}5°$};
  % echo primes
  \fill[gray] ({cos(21.9)},{sin(21.9)}) circle (0.4pt);
  \draw[gray, thin] ({cos(21.9)},{sin(21.9)}) -- (0.58, 0.22);
  \node[right, font=\scriptsize, gray] at (0.58, 0.22) {$p\!=\!11$ (écho)};
  \fill[gray] ({cos(20.8)},{sin(20.8)}) circle (0.4pt);
  \draw[gray, thin] ({cos(20.8)},{sin(20.8)}) -- (0.58, 0.10);
  \node[right, font=\scriptsize, gray] at (0.58, 0.10) {$p\!=\!13$ (écho)};
  % annotation
  \node[font=\small, align=center] at (0.05, 1.08)
    {$\cos\theta_p = 1 - \delta_p$};
\end{tikzpicture}
\caption{Angles d'holonomie sur le cercle unité.
  Chaque premier actif $p \in \{3,5,7\}$ définit un angle $\theta_p$
  via $\cos\theta_p = 1 - \delta_p$.
  Le sinus au carré $\sin^2\!\theta_p$ est la force de couplage.
  Les premiers d'écho ($p \geq 11$, gris) sont plus proches de
  l'axe.}
\label{fig:coupling_circle}
\end{figure}

% --- Figure: active prime threshold ---
\begin{figure}[htbp]
\centering
\begin{tikzpicture}[xscale=2.0, yscale=4.5, >=stealth]
  % axes
  \draw[->,thick] (0,0) -- (6.5,0) node[right,font=\small] {premier $p$};
  \draw[->,thick] (0,0) -- (0,0.95) node[above,font=\small] {$\gamma_p$};
  % threshold line s = 1/2
  \draw[condorange, thick, dashed] (0,0.5) -- (6.3,0.5)
    node[right, font=\small] {$s = 1/2$};
  % bars
  \fill[thmblue!70] (0.7,0) rectangle (1.3,0.808);
  \node[above, font=\small, thmblue] at (1,0.808) {$0{,}808$};
  \node[below, font=\small] at (1,-0.02) {$3$};
  \fill[thmblue!70] (1.7,0) rectangle (2.3,0.696);
  \node[above, font=\small, thmblue] at (2,0.696) {$0{,}696$};
  \node[below, font=\small] at (2,-0.02) {$5$};
  \fill[thmblue!70] (2.7,0) rectangle (3.3,0.595);
  \node[above, font=\small, thmblue] at (3,0.595) {$0{,}595$};
  \node[below, font=\small] at (3,-0.02) {$7$};
  \fill[gray!50] (3.7,0) rectangle (4.3,0.426);
  \node[above, font=\small, gray] at (4,0.426) {$0{,}426$};
  \node[below, font=\small] at (4,-0.02) {$11$};
  \fill[gray!50] (4.7,0) rectangle (5.3,0.356);
  \node[above, font=\small, gray] at (5,0.356) {$0{,}356$};
  \node[below, font=\small] at (5,-0.02) {$13$};
  % labels
  \node[font=\small, thmblue] at (2,0.88) {\textbf{Actif} ($\gamma_p > s$)};
  \node[font=\small, gray] at (4.5,0.88) {\textbf{Écho} ($\gamma_p < s$)};
  % y-axis ticks
  \foreach \y in {0.2, 0.4, 0.6, 0.8} {
    \draw (-0.05, \y) -- (0.05, \y);
    \node[left, font=\scriptsize] at (-0.07, \y) {$\y$};
  }
\end{tikzpicture}
\caption{Dimensions anomales $\gamma_p$ à $\mustar = 15$.
  Le seuil $\gamma_p = s = 1/2$ (tirets orange) sépare les trois
  premiers actifs $\{3,5,7\}$ (bleu) des premiers d'écho
  $\{11,13,\ldots\}$ (gris). Ceci détermine $N_c = 3$.}
\label{fig:active_threshold}
\end{figure}

% --- Figure: 3D coupling cube ---
\begin{figure}[htbp]
\centering
\tdplotsetmaincoords{70}{120}
\begin{tikzpicture}[tdplot_main_coords, scale=4.5]
  % Unit sphere (wireframe) — domain [0,1]^3 boundary
  % Equatorial circle in xy-plane (z=0.5)
  \draw[gray!25, thin] plot[domain=0:360, samples=72, smooth cycle]
    ({0.5 + 0.5*cos(\x)}, {0.5 + 0.5*sin(\x)}, 0.5);
  % Meridian circle in xz-plane (y=0.5)
  \draw[gray!25, thin] plot[domain=0:360, samples=72, smooth cycle]
    ({0.5 + 0.5*cos(\x)}, 0.5, {0.5 + 0.5*sin(\x)});
  % Meridian circle in yz-plane (x=0.5)
  \draw[gray!25, thin] plot[domain=0:360, samples=72, smooth cycle]
    (0.5, {0.5 + 0.5*cos(\x)}, {0.5 + 0.5*sin(\x)});
  % Axes
  \draw[->, thick, predred] (0,0,0) -- (1.25,0,0)
    node[below, font=\small] {$\sin^2\!\theta_3$};
  \draw[->, thick, thmblue] (0,0,0) -- (0,1.25,0)
    node[right, font=\small] {$\sin^2\!\theta_5$};
  \draw[->, thick, bridgepurple] (0,0,0) -- (0,0,1.25)
    node[above, font=\small] {$\sin^2\!\theta_7$};
  % The alpha_EM point: exact values from delta_p(2-delta_p) at q=13/15, mu*=15
  \pgfmathsetmacro{\sa}{0.2192}
  \pgfmathsetmacro{\sb}{0.1940}
  \pgfmathsetmacro{\sc}{0.1726}
  % Three circles intersecting at alpha point (radius 0.3)
  % Circle in sin^2(theta_3) = const plane (yz-plane at x=sa)
  \draw[predred, thick, fill=predred!8, fill opacity=0.3]
    plot[domain=0:360, samples=36, smooth cycle]
    (\sa, {\sb + 0.3*cos(\x)}, {\sc + 0.3*sin(\x)});
  % Circle in sin^2(theta_5) = const plane (xz-plane at y=sb)
  \draw[thmblue, thick, fill=thmblue!8, fill opacity=0.3]
    plot[domain=0:360, samples=36, smooth cycle]
    ({\sa + 0.3*cos(\x)}, \sb, {\sc + 0.3*sin(\x)});
  % Circle in sin^2(theta_7) = const plane (xy-plane at z=sc)
  \draw[bridgepurple, thick, fill=bridgepurple!8, fill opacity=0.3]
    plot[domain=0:360, samples=36, smooth cycle]
    ({\sa + 0.3*cos(\x)}, {\sb + 0.3*sin(\x)}, \sc);
  % Projection lines to axes
  \draw[predred, densely dashed, thin] (\sa,\sb,\sc) -- (\sa,0,0);
  \draw[thmblue, densely dashed, thin] (\sa,\sb,\sc) -- (0,\sb,0);
  \draw[bridgepurple, densely dashed, thin] (\sa,\sb,\sc) -- (0,0,\sc);
  % The coupling point
  \fill[condorange] (\sa,\sb,\sc) circle (1.0pt);
  % Leader line to label
  \draw[condorange, thin] (\sa,\sb,\sc) -- (0.55,0.55,0.65);
  \node[right, font=\small, condorange, align=left] at (0.55,0.55,0.65)
    {$\alphaEM = \prod \sinp{p}$\\[1pt]\footnotesize$= 1/137{,}036$};
  % Axis tick labels
  \node[below, font=\tiny, predred] at (\sa,0,0) {$0{,}219$};
  \node[left, font=\tiny, thmblue] at (0,\sb,0) {$0{,}194$};
  \node[left, font=\tiny, bridgepurple] at (0,0,\sc) {$0{,}172$};
  % Origin label
  \node[below left, font=\tiny, gray] at (0,0,0) {$0$};
\end{tikzpicture}
\caption{La sphère de couplage.
  Chaque axe représente $\sinp{p}$ pour un premier actif.
  Trois cercles colorés, chacun perpendiculaire à un axe, se coupent
  au point unique $(\sinp{3}, \sinp{5}, \sinp{7})$ dont le produit
  vaut $\alphaEM = 1/137{,}036$ (\BRIDGETAG, voir~\cite{companion_M5}).
  Le point de couplage se trouve profondément dans la région des
  petits $\theta$~: la nature n'utilise qu'une faible fraction de
  l'espace de paramètres disponible.}
\label{fig:coupling_cube}
\end{figure}

% ============================================================
\subsection{$N_c = 3$~: nombre de premiers actifs}
\label{sec:Nc_3}

\mTHM
\begin{theoreme}[$N_c = 3$]
\label{thm:Nc_3}\leanProved{PT.FixedPoint.T7MuStar.activeAt15\_eq}
Le nombre de premiers actifs à $\mustar = 15$ est $N_c = 3$.

\begin{proof}
Par le théorème~\ref{thm:active}~: exactement trois premiers
satisfont $\gamma_p > s = 1/2$ à $\mustar = 15$, à savoir
$\{3, 5, 7\}$. Tous les premiers $p \geq 11$ ont
$\gamma_p < 1/2$. Donc
$N_c := |\{p : \gamma_p(\mustar) > s\}| = 3$.
\end{proof}
\end{theoreme}

\begin{corollaire}[Coefficient de Casimir]
\label{cor:CF}
$C_F = (N_c^2 - 1)/(2N_c) = 4/3$.
C'est une \emph{conséquence} de $N_c = 3$, non une entrée.
\end{corollaire}

\begin{remarque}[Contrôle algébrique de cohérence (R14)]
\label{rem:Nc_algebraic}
Le résultat peut être recoupé algébriquement sans supposer $C_F$.
T1 interdit 2 auto-transitions à $p = 3$ (classes $1$ et $2$),
laissant $n_{\rm perm} = N_c^2 - 2 = 7$ transitions permises sur~9.
En posant $n_{\rm perm}/N_c = C_F + 1$ avec
$C_F = (N_c^2 - 1)/(2N_c)$~:
\[
  \frac{N_c^2 - 2}{N_c} = \frac{N_c^2 + 2N_c - 1}{2N_c},
\]
ce qui se simplifie en $N_c^2 - 2N_c - 3 = 0$, c'est-à-dire
\begin{equation}
  (N_c - 3)(N_c + 1) = 0 \implies N_c = 3 \quad (N_c > 0).
  \label{eq:Nc_equation}
\end{equation}
La preuve par comptage et l'équation algébrique donnent toutes deux
$N_c = 3$ indépendamment. (Vérification~:
\texttt{ch06\_holonomy/test\_D07\_sin2\_identity.py}.)
\end{remarque}

\begin{remarque}[Cinq voies indépendantes vers $N_c = 3$]
\label{rem:Nc_routes}
Au-delà du comptage des premiers actifs et du contrôle algébrique
ci-dessus, trois voies supplémentaires confirment $N_c = 3$ (R38b)~:
\begin{enumerate}[nosep]
\item $4N_c + 3 = \mustar$~: comptage des fermions de Weyl par
  génération (BRIDGE).
\item $(N_c + 2)(N_c + 3) = 2\mustar = 30$~: décomposition SU(5)
  $\bar{5} + 10 = 15$ (BRIDGE).
\item $(N_c + 1)!/(N_c + 3) = 2^{N_s - 1}$~: seul $N_c = 3$ donne
  une puissance entière de~2, fournissant $N_s = 3$ dimensions
  spatiales (BRIDGE).
\end{enumerate}
Le comptage des premiers actifs ($|\{p : \gamma_p > s\}| = 3$,
\THMTAG) et le contrôle algébrique ($N_c = 3$ depuis l'équation,
\THMTAG) sont des résultats arithmétiques. Les trois voies
supplémentaires (items~1--3) impliquent des identifications
physiques (\BRIDGETAG).
\end{remarque}

\begin{remarque}[Identification physique]
L'identification de $N_c = 3$ (dérivé de la structure du crible) avec
le nombre de couleurs QCD est un énoncé de pont
(\cite{companion_M5}). Ici nous ne prouvons que le résultat
arithmétique selon lequel $N_c = 3$.
\end{remarque}

% ============================================================
\subsection{Premiers d'écho et corrections perturbatives}
\label{sec:ghost}

Les premiers inactifs $\{11, 13, \ldots\}$ ne sont pas simplement
absents. Ils apparaissent comme \emph{corrections perturbatives}
(polarisation du vide d'écho).

\begin{definition}[Poids VP d'écho]
\label{def:ghost_vp}
Le poids de polarisation du vide d'écho est~:
\begin{equation}
  \beta_{\rm echo} = \sum_{p \in \{11,13\}} \sinsp{p}(\qplus) \cdot \gamma_p
  \approx 0{,}104.
  \label{eq:ghost_vp:ch06}
\end{equation}
\end{definition}

La contribution VP d'écho à $\alpha_{\rm EM}$ est
$\delta\alpha/\alpha = -\gamma_3 \cdot \alpha_{\rm bare} \cdot \beta_{\rm echo} \approx 14\;{\rm ppb}$
(voir~\cite{companion_physics}).

% ============================================================
\subsection{L'origine arithmétique de $\pi$}
\label{sec:pi_origin}

L'identité d'holonomie révèle un fait structurel à propos de $\pi$~:
la constante $\pi$ n'est pas importée dans le crible depuis la
géométrie euclidienne. Elle \emph{émerge} de l'arithmétique de la
divisibilité.

\begin{remarque}[Pourquoi le cercle apparaît]
\label{rem:why_circle}
La matrice de transfert $T_p$ agit sur $\mathbb{Z}/p\mathbb{Z}$ en
redistribuant la masse de probabilité entre classes de résidus. Sa
restriction $2 \times 2$ aux classes non nulles (après que T1 a
forcé la diagonale à zéro à $p = 3$) a des entrées diagonales
$\cos\theta_p = 1 - \delta_p$ et des entrées hors-diagonales
$\sin^2\!\theta_p = \delta_p(2 - \delta_p)$.
Celles-ci satisfont l'identité de Pythagore
$\cos^2\theta + \sin^2\theta = 1$ \emph{automatiquement}~: c'est une
conséquence de la contrainte de stochasticité $\sum_j T_{ij} = 1$
sur la matrice de transfert, et non une hypothèse sur les cercles.

Le cercle unité $S^1$ est donc la \emph{variété de contraintes} des
amplitudes de transition du crible. Sa période $2\pi$ est déterminée
par l'exigence qu'une rotation complète ramène à l'identité. Le
crible n'utilise pas $\pi$~--- il \emph{engendre} $S^1$, et $\pi$
est la demi-période de ce cercle.
\end{remarque}

\begin{remarque}[Calculer $\pi$ depuis le crible]
\label{rem:pi_computation}
Le produit eulérien pour $\zeta(2)$ admet une interprétation directe
en termes de crible. À chaque premier~$p$, le crible élimine les
multiples de~$p$~; la probabilité de survie \emph{au carré} à travers
tous les premiers donne la densité des paires premières avec chaque
premier~:
\begin{equation}
  \prod_{p\;\text{premier}} \frac{1}{1 - 1/p^2} = \zeta(2) = \frac{\pi^2}{6}.
  \label{eq:euler_pi}
\end{equation}
Chaque facteur $(1 - 1/p^2)^{-1}$ est un rapport de résidus du
crible~: la probabilité que deux entiers indépendants survivent
tous deux au filtre en~$p$, divisée par la base non contrainte. Le
produit infini converge vers $\pi^2/6$~--- un énoncé que l'effet
cumulé de tous les filtres arithmétiques, composés
multiplicativement, encode la géométrie de $S^1$.

En tronquant le produit au $k$-ième premier $p_k$, on obtient une
\emph{approximation par crible} de $\pi$~:
\begin{equation}
  \pi_k = \sqrt{6 \prod_{p \leq p_k} \frac{1}{1 - 1/p^2}},
  \qquad \pi_k \to \pi \;\;\text{lorsque } k \to \infty.
  \label{eq:pi_sieve}
\end{equation}
La convergence est lente ($|\pi_k - \pi| \sim 3/p_k$), mais le
point conceptuel est net~: $\pi$ est la valeur limite d'un produit
purement arithmétique sur les premiers, sans entrée géométrique.
Vérification~: \texttt{ch09\_bridge/test\_equations\_pont\_PT.py},
test~H5 (24 premiers, $|\pi_{24} - \pi|/\pi < 0{,}5\%$).
\end{remarque}

\begin{remarque}[Pourquoi $\pi$ apparaît dans les formules sur les premiers]
\label{rem:why_pi_in_primes}
La surprise classique~--- pourquoi $\pi$ apparaît-il dans
$\zeta(2)$, dans le théorème de Mertens, dans la série singulière
de Hardy--Littlewood, dans le terme d'erreur du théorème des nombres
premiers~?~--- se dissout dans le cadre de l'holonomie.

À chaque premier $p$, le crible impose une \emph{rotation}
$\theta_p$ sur l'espace des résidus. L'holonomie cumulée de ces
rotations est mesurée par des produits de fonctions
trigonométriques de $\theta_p$. Puisque $\sin$ et $\cos$ sont
périodiques de période $2\pi$, toute formule multiplicative sur les
premiers hérite de~$\pi$ comme \emph{période de la rotation que la
divisibilité engendre}.

En particulier~:
\begin{itemize}[nosep]
\item $\zeta(2) = \pi^2/6$~: produit cumulé de survie au carré
  (\cref{eq:euler_pi}).
\item $\prod_{p \leq x}(1 - 1/p) \sim 2e^{-\gamma}/\ln x$
  (Mertens)~: la fraction de survie décroît logarithmiquement, et
  la constante $e^{-\gamma}$ encode la même accumulation
  d'holonomie.
\item $\alphaEM = \prod_{p \in \{3,5,7\}} \sinp{p}$~: la constante
  de structure fine est un produit de valeurs $\sin^2\!\theta_p$~---
  littéralement un produit d'amplitudes d'holonomie sur~$S^1$.
\item $\mu_{\rm end} = N_c \cdot \pi = 3\pi$~: la borne
  d'intégration RG vaut $N_c$ demi-tours (phase de Berry,
  voir~\cite{companion_physics}).
\end{itemize}

Dans chaque cas, $\pi$ entre en théorie des nombres non comme un
import depuis la géométrie, mais comme la \emph{période intrinsèque
de la dynamique rotationnelle que le crible crée sur les espaces de
résidus}. L'arithmétique contraint la géométrie~; la géométrie
restitue $\pi$.
\end{remarque}

% ============================================================
\subsection{Synthèse et connexions}
\label{sec:ch6_summary}

\begin{ledgerbox}
Noyau dur~:
  Identité d'holonomie~: $\sinp{p} = \delta_p(2-\delta_p)$~---
  identité algébrique.
  Critère de premier actif~: $\{3,5,7\}$ actifs ($\gamma_p > s$)~---
  THM numérique à $\mustar=15$.
  $N_c = 3$~: comptage des premiers actifs à $\mustar = 15$~--- THM
  arithmétique.
  Contrôle algébrique~: $(N_c-3)(N_c+1)=0$ depuis le comptage des
  transitions T1.

Conditionnel~:
  Les valeurs numériques de $\sinp{p}$ et $\gamma_p$ utilisent
  $\mustar = 15$ (\cite{companion_M3}, T5 par balayage).

Identification (THM, quatre unicités)~:
  $\sinp{p}$ = constante de couplage (\cite{companion_M5}, BA3,
  fermée par G1+G3).
  $N_c = 3$ = couleurs QCD (\cite{companion_M5}, fermée par T5+T6).
  $\{11,13\}$ = contre-termes VP d'écho (\cite{companion_physics},
  DER-PHYS).

Validation~:
  $\sinp{p}$ vérifié par
  \texttt{ch06\_holonomy/test\_D07\_sin2\_identity.py}.
  $N_c = 3$ vérifié par
  \texttt{ch06\_holonomy/test\_D07\_sin2\_identity.py}.
  VP d'écho 14\,ppb vérifié (\cite{companion_physics}).
\end{ledgerbox}

\medskip\noindent\textbf{Connexions vers l'aval.}
\begin{itemize}
\item \cite{companion_M3} fournit $\mustar = 15$ (nécessaire pour
  les valeurs numériques ici).
\item \cite{companion_M5} dérive $\sinsp{p}$ comme constantes de
  couplage (DER, Pontryagin + Wightman).
\item L'article compagnon~\cite{companion_physics} assemble le
  produit $\prod_{p=3,5,7} \sinsp{p}$ pour dériver
  $\alpha_{\rm EM}$.
\item L'article compagnon~\cite{companion_physics} utilise
  $\gamma_p$ pour dériver $\sinW$ et $\alpha_s$.
\item L'article compagnon~\cite{companion_physics} utilise les
  directions actives pour définir la métrique spatiale.
\end{itemize}


\newpage

\section{Conclusion}
\label{sec:conclusion}

Cet article a établi deux piliers de la théorie de la persistance qui
relient la structure combinatoire du crible à des objets géométriques
et spectraux.

\textbf{Géométrie de l'information.}
La structure produit TCR de l'espace d'états du crible $\Delta_m$
sélectionne uniquement deux objets géométriques~: la divergence de KL
$\DKL$ (théorème~G1, par réduction à Shore--Johnson et preuve autonome
via TCR) et la métrique de Fisher $\Fisher$ (théorème~G3, par
réduction à \v{C}encov). Le lien entre les deux est classique
(théorème~G2~: Fisher est la hessienne de $\DKL$). La contribution
clé de PT n'est pas de prouver Shore--Johnson ou \v{C}encov~--- ce
sont des théorèmes classiques importés~--- mais de vérifier que
l'espace d'états du crible, avec sa factorisation TCR et ses matrices
de transfert de Markov, satisfait les prémisses des deux théorèmes
d'unicité. Sept tests d'élimination confirment qu'aucune divergence ou
métrique alternative n'est compatible avec les symétries du crible.

La métrique de Fisher encode la dualité entre addition et
multiplication~: c'est l'unique géométrie qui respecte les deux
opérations simultanément. L'identité du TFS $\log_2 m = \DKL + H$ est
cette dualité exprimée comme loi de conservation~; la géométrie de
Fisher en est l'incarnation comme métrique riemannienne.

\textbf{Holonomie.}
L'identité d'holonomie $\sinsp{p} = \delta_p(2 - \delta_p)$ est une
tautologie algébrique qui reçoit trois dérivations indépendantes~:
géométrique (identité de Pythagore depuis la stochasticité), spectrale
(contraction de Fourier du caractère fondamental) et
information-géométrique (composante de courbure de Fisher par
premier). L'accord des trois voies élimine le risque de preuve unique.

Les dimensions anomales $\gamma_p$ sont strictement décroissantes en
$p$, et le critère de premier actif $\gamma_p > s = 1/2$ sélectionne
exactement $\{3, 5, 7\}$ au point fixe $\mustar = 15$
(voir~\cite{companion_M3}). Ceci donne $N_c = 3$ comme théorème, et
non comme entrée. Le coefficient de Casimir $C_F = 4/3$ en découle
comme corollaire. Un contrôle algébrique indépendant~---
$(N_c - 3)(N_c + 1) = 0$ depuis le comptage des transitions T1~---
confirme le résultat sans référence aux dimensions anomales.

Les trois premiers actifs définissent trois directions d'holonomie
indépendantes. Leur identification physique comme dimensions
spatiales, et l'identification de $N_c = 3$ avec le nombre de
couleurs QCD, sont des énoncés de pont développés
dans~\cite{companion_M5}.

\textbf{Questions ouvertes.}
Le critère de seuil $\gamma_p > s$ est naturel (cinq arguments
indépendants, intervalle de stabilité large de $35\%$) mais manque
d'une preuve d'optimalité sous forme close. Une dérivation rigoureuse
montrant que $\tau = s$ est le \emph{seul} seuil cohérent~--- plutôt
que simplement le choix le plus naturel dans une large bande de
stabilité~--- reste ouverte.

Les premiers d'écho ($p \geq 11$) contribuent perturbativement comme
corrections de polarisation du vide. La question de savoir si la
série VP d'écho converge absolument, et si elle admet une
resommation sous forme close, est un travail futur.

\textbf{Vers l'aval.}
L'article compagnon~\cite{companion_M3} prouve la convergence des
récurrences du crible et établit le point fixe $\mustar = 15$ qui
sous-tend les valeurs numériques utilisées tout au long de cet
article. Les structures mathématiques étendues~--- mécanique
complexe, graphes spectraux et extensions de symétrie~--- sont
développées dans~\cite{companion_M4}. Le pont de l'arithmétique vers
la physique, où les objets géométriques établis ici acquièrent un
sens physique, est l'objet de~\cite{companion_M5}.

\bibliographystyle{plainnat}
\bibliography{references}

\end{document}
